\section{Methodology}
\label{sec:methodology}

The method is a two-stage tabular prediction workflow rather than a new base learning algorithm. Its main design choice is to separate an atmospheric-transmissivity estimate from the final plant-power estimate. This keeps the model close to the physical structure of PV generation: sun position and plant rating define a resource envelope, while meteorological and cloud conditions explain how much of that envelope is available at a given timestamp.

\subsection{Forecasting Problem and Feature Availability}

For each timestamp $t$ on a 5-min target grid, the proposed path predicts AC plant power $\widehat{\Ppv}(t)$ using calendar variables, simplified solar-position variables, meteorological variables, and cloud-layer variables available for that target timestamp. In the offline test reported here, these target-time weather and cloud columns are taken from the held-out dataset to measure model behavior under a fixed feature-engineering protocol. In deployment, the same fields must be supplied before issue time by a forecast service or local weather pipeline; no archived issue-time weather forecasts or lead-time-specific forecast errors are evaluated in this manuscript. The proposed path does not use measured irradiance or measured plant power at timestamp $t$. A separate sensor-assisted model is reported only as an upper-bound diagnostic because it uses measured irradiance-derived features at the target timestamp.

\begin{table*}[!t]
\caption{Feature availability by evaluation path. The proposed path is the paper's deployable forecast path and assumes meteorological and cloud inputs for the target timestamp are supplied by a forecast service or local weather pipeline.}
\label{tab:feature_availability}
\centering
\scriptsize
\setlength{\tabcolsep}{3pt}
\begin{tabularx}{\textwidth}{@{}>{\raggedright\arraybackslash}p{0.24\textwidth}*{4}{>{\centering\arraybackslash}X}@{}}
\toprule
Input group & Proposed forecast & Sensor-assisted diagnostic & 5-min persistence & Previous-day persistence \\
\midrule
Calendar/time features & yes & yes & no & time index only \\
Simplified solar geometry & yes & yes & no & no \\
Meteorological inputs & yes & yes & no & no \\
Cloud total/low/mid/high & yes & yes & no & no \\
Stage-1 estimated $\kt$ and irradiance proxy & yes & no & no & no \\
Raw logged \texttt{theoretical\_power} & no & yes & no & no \\
Measured irradiance at target time & no & yes & no & no \\
Measured irradiance lags/rolling windows & no & yes & no & no \\
Live plant power at $t-5$ min & no & no & yes & no \\
Plant power from $t-24$ h & no & no & no & yes \\
\bottomrule
\end{tabularx}
\end{table*}


Table~\ref{tab:feature_availability} explains why the baselines are not interchangeable. The 5-min last-value persistence baseline uses live plant power at $t-5$ min, so it is a nowcast baseline with information that the proposed target-weather-input path intentionally does not use. The previous-day same-time persistence baseline uses $P(t-24~\mathrm{h})$, so it is closer to a simple schedule-oriented baseline.

\begin{figure*}[!t]
\centering
\includegraphics[width=0.88\textwidth]{figures/diagrams/diagram_workflow.png}
\caption{Data and modeling workflow used to construct the two-stage forecasting pipeline. The same feature logic is used for training diagnostics and forecast-CSV inference.}
\label{fig:workflow}
\end{figure*}

\begin{figure}[!t]
\centering
\includegraphics[width=\columnwidth]{figures/diagrams/feature_availability_matrix.png}
\caption{Feature availability matrix for the proposed forecast path, the sensor-assisted diagnostic path, and the two persistence baselines.}
\label{fig:feature_availability_matrix}
\end{figure}

\subsection{Solar Geometry and Clear-Sky Reference}

The timestamp convention is Indian Standard Time (IST, UTC+05:30). The implementation does not apply an equation-of-time correction or a longitude correction to convert civil clock time to local apparent solar time. This simplified convention is retained for reproducibility, and its consequences are listed in Section~\ref{sec:limitations}. The implementation evaluates the solar-position equations in radians as follows.

For day $n$ of the year, solar declination is approximated as
\begin{equation}
\delta_n^{\mathrm{rad}} =
\frac{\pi}{180}23.45
\sin\left(\frac{2\pi}{365}(n-81)\right).
\label{eq:declination}
\end{equation}
For civil hour $h(t)$, including minutes as fractional hours, the hour angle is
\begin{equation}
\omega^{\mathrm{rad}}(t)=\frac{\pi}{180}15\left(h(t)-12\right).
\label{eq:hour_angle}
\end{equation}
Given latitude $\phi$ in radians, the intermediate elevation sine is
\begin{equation}
s_{\alpha}(t)=
\sin\phi\sin\delta_n^{\mathrm{rad}}
+\cos\phi\cos\delta_n^{\mathrm{rad}}\cos\omega^{\mathrm{rad}}(t),
\label{eq:elevation}
\end{equation}
and the solar elevation angle used for thresholding is
\begin{equation}
\alpha(t)=\arcsin\left(\mathrm{clip}\left(s_{\alpha}(t),-1,1\right)\right).
\label{eq:elevation_angle}
\end{equation}
The clear-sky irradiance proxy is
\begin{equation}
\Gcs(t)=1000\max\left(0,s_{\alpha}(t)\right),
\label{eq:gcs}
\end{equation}
which is naturally nonnegative and bounded by 1000 W/m$^2$ after the sine clipping used in \eqref{eq:elevation_angle}. The corresponding clear-sky power feature is
\begin{equation}
P_{\mathrm{cs}}(t)=\frac{A_{\mathrm{eff}}\Gcs(t)}{1000},
\label{eq:pcs}
\end{equation}
where $A_{\mathrm{eff}}=181.591$ kW per kW/m$^2$ is the effective irradiance-to-power coefficient stored in the active model metadata. It is a calibrated envelope coefficient used by the feature pipeline, not a claim about the physical module area alone. The air-mass-like feature used by the tabular model is
\begin{equation}
M(t)=
\begin{cases}
1/\mathrm{clip}\left(s_{\alpha}(t),0.05,1\right), & s_{\alpha}(t)>0.05,\\
20, & s_{\alpha}(t)\leq0.05.
\end{cases}
\label{eq:airmass}
\end{equation}
These terms are reproducible physics proxies, not a full pvlib clear-sky model. The raw dataset column \texttt{theoretical\_power} is kept as a logged reference channel and is used in the sensor-assisted diagnostic path; the proposed path uses the deterministic $P_{\mathrm{cs}}(t)$ feature and the derived Stage-1 envelope $\widehat{P}_{\mathrm{theo},c}(t)$ instead of using raw \texttt{theoretical\_power} as a target-time input.

\subsection{Stage-1 Transmissivity Target}

The Stage-1 target is a clearness-index-like transmissivity:
\begin{equation}
\kt(t)=
\begin{cases}
\mathrm{clip}\left(\dfrac{G(t)}{\Gcs(t)},0,1.2\right), & \Gcs(t)>10,\\
0, & \Gcs(t)\leq 10,
\end{cases}
\label{eq:kt}
\end{equation}
where $G(t)$ is the dataset irradiance channel. Stage 1 is trained only on normal daylight rows above the solar-elevation threshold stored in the metadata; this threshold is expressed in degrees after evaluating \eqref{eq:elevation_angle}. At inference time, Stage 1 estimates $\widehat{\kt}(t)$ from meteorological, time, cloud-layer, and solar-geometry features:
\begin{equation}
\widehat{\kt}(t)=f_1\left(\mathbf{x}_{\mathrm{weather}}(t),
\mathbf{x}_{\mathrm{time}}(t),\mathbf{x}_{\mathrm{solar}}(t),
\mathbf{x}_{\mathrm{cloud}}(t)\right).
\label{eq:stage1}
\end{equation}
Because the Stage-1 regressor itself is not range constrained, its output is clipped before constructing the power envelope:
\begin{equation}
\widehat{\kt}_{c}(t)=\mathrm{clip}\left(\widehat{\kt}(t),0,1.2\right).
\label{eq:kt_clip}
\end{equation}
The estimated irradiance proxy is then
\begin{equation}
\widehat{G}_{c}(t)=\mathrm{clip}\left(\Gcs(t)\widehat{\kt}_{c}(t),0,1100\right),
\label{eq:irr_est}
\end{equation}
and the estimated theoretical power envelope is
\begin{equation}
\widehat{P}_{\mathrm{theo},c}(t)=
\max\left(0,\frac{A_{\mathrm{eff}}\widehat{G}_{c}(t)}{1000}\right).
\label{eq:theo_est}
\end{equation}

\subsection{Lag, Rolling, and Weather-Derived Features}

For a variable $x(t)$ sampled at $\Delta t=5$ min, lag features are
\begin{equation}
x_{\ell}(t)=x(t-\ell\Delta t),
\label{eq:lag}
\end{equation}
rolling means over a window of $m$ samples are
\begin{equation}
\bar{x}_{m}(t)=\frac{1}{m}\sum_{i=0}^{m-1}x(t-i\Delta t).
\label{eq:rolling}
\end{equation}
Accumulated-window features are defined separately as
\begin{equation}
S^{x}_{m}(t)=\sum_{i=0}^{m-1}x(t-i\Delta t).
\label{eq:rolling_sum}
\end{equation}
The proposed path includes 1-h temperature and humidity lags, 1-h/3-h/24-h precipitation sums using \eqref{eq:rolling_sum}, 3-h humidity statistics using \eqref{eq:rolling}, total cloud-cover lag and trend, and first and second harmonic calendar encodings. It excludes measured plant-power lags and measured-irradiance lags or rolling features. Those measured resource features are confined to the sensor-assisted diagnostic path, while the proposed path uses only deterministic solar proxies and Stage-1 estimated irradiance features. Dew point is approximated as
\begin{equation}
T_{\mathrm{dew}}(t)=T(t)-\frac{100-\mathrm{RH}(t)}{5}.
\label{eq:dewpoint}
\end{equation}
The simplified dew-point equation is retained for consistency with model training; replacing it with a Magnus formula would define a different feature set.

\subsection{Stage-2 Power Model and Post-Processing}

Stage 2 estimates plant power from meteorological features, solar features, and the Stage-1 estimated atmospheric variables:
\begin{multline}
\widetilde{\Ppv}(t)=f_2\bigl(\widehat{\kt}_{c}(t),\widehat{G}_{c}(t),
\widehat{P}_{\mathrm{theo},c}(t),\\
\mathbf{x}_{\mathrm{weather}}(t),\mathbf{x}_{\mathrm{solar}}(t),
\mathbf{x}_{\mathrm{lag}}(t)\bigr).
\label{eq:stage2}
\end{multline}
For the proposed path, $\mathbf{x}_{\mathrm{lag}}(t)$ contains only lagged or rolling meteorological variables, cloud variables, deterministic solar proxies, and Stage-1 estimated irradiance features available from the target-weather-input pipeline. It does not contain measured plant power, measured target-time irradiance, or measured irradiance rolling windows. The sensor-assisted diagnostic model uses measured irradiance-derived rolling features separately and is therefore not used for the headline deployable-path claim.
The final reported prediction is explicitly post-processed as
\begin{equation}
\widehat{\Ppv}(t)=\min\left\{\Prated,\widehat{P}_{\mathrm{theo},c}(t),
\max\left(0,\widetilde{\Ppv}(t)\right)\right\}.
\label{eq:postprocess}
\end{equation}
All headline metrics in this paper are computed after Stage-1 clipping, envelope construction, and the final clipping operation in \eqref{eq:postprocess}. The paper therefore does not claim that unconstrained XGBoost regression guarantees bounded predictions; boundedness is supplied by the explicit post-processing rules in \eqref{eq:kt_clip}--\eqref{eq:postprocess}.

\subsection{XGBoost Objective}

Both stages use XGBoost regression. At boosting iteration $m$, the regularized objective is
\begin{equation}
\mathcal{L}^{(m)} =
\sum_{i=1}^{N} l\left(y_i,\hat{y}_i^{(m-1)}+f_m(\mathbf{x}_i)\right)
+ \Omega(f_m),
\label{eq:xgb_obj}
\end{equation}
where squared-error loss is used for the regression task. The tree penalty is
\begin{equation}
\Omega(f_m)=\gamma T_m+\frac{1}{2}\lambda\sum_{j=1}^{T_m} w_j^2,
\label{eq:xgb_reg}
\end{equation}
where $T_m$ is the number of leaves and $w_j$ are leaf weights. The regularization controls tree complexity in a dataset with correlated weather, time, and solar-position features.

\subsection{Algorithm Summary}

The operational pipeline is:
\begin{enumerate}
\item sort and audit the 5-min dataset by timestamp;
\item compute time, solar-geometry, weather-anomaly, lag, rolling, and cloud-derived features;
\item split chronologically into training, validation, and held-out test blocks;
\item train Stage 1 to estimate $\kt$ from meteorological, time, solar, and cloud features;
\item apply Stage-1 clipping, then form estimated irradiance and theoretical-power features;
\item train Stage 2 to predict plant power from the estimated solar features and meteorological features;
\item clip the output using \eqref{eq:postprocess};
\item evaluate on common daylight masks against persistence and ablation baselines.
\end{enumerate}
